% GENERATED FILE. Edit canonical lessons, then run npm run generate:books.
% canonical-source-hash: fnv1a64:a908a30473f81d21
% canonical-lessons: ZH-C519-chunguang, ZH-C519-qingchen, ZH-C519-sijiruchun, ZH-C519-xuejing, ZH-C519-chengliang, ZH-R519-review-514-519

\chapter{Out in Nature}
\label{ch:zh-out-in-nature}
\hlchaptermodality{519}

\begin{chapteropening}
\textbf{By the end of this chapter:} \emph{I can talk about spring at its best, the early morning, a place that is mild all year, a snowy view and sitting in the cool.}

Complete the last lesson of chapter 519: I can talk about spring at its best, the early morning, a place that is mild all year, a snowy view and sitting in the cool.
\end{chapteropening}

\section[\texorpdfstring{\zh{春光} (chūnguāng)}{春光 (chūnguāng)}]{\zh{春光} (chūnguāng) — the beauty of spring; spring sunshine}

\label{lesson:ZH-C519-chunguang}



\begin{quote}
Before the new one: say the Mandarin for to put up an umbrella, then the Mandarin for a thick fog.
\end{quote}

\subsection*{You'll want to know: \zh{春光}}
\textbf{\zh{春光}} — \emph{chūnguāng} — \textquotedblleft{}the beauty of spring; spring sunshine\textquotedblright{}.

\begin{quote}
\textbf{\zh{春光最好的时候树都开花了}} \emph{chūnguāng zuì hǎo de shíhou shù dōu kāihuā le} — When spring is at its loveliest, the trees are all in blossom.
\end{quote}

\begin{cousinweb}
\textbf{\zh{春光}} (\emph{chūnguāng}) is built from \emph{chūn} 'spring' and \zh{光} 'light', both inherited from Middle Chinese, \emph{tsyhwin} and \emph{kwang} (Guangyun, in Baxter's transcription), which glosses the first, after an older book, as the head of the four seasons, when all things come out.
\end{cousinweb}

\begin{grammarlens}[title={What you've built}]
One more word for saying how it was.
\end{grammarlens}

\subsection*{Your turn}
\begin{itemize}
\raggedright
  \item \emph{Say it:} \emph{chūnguāng}
  \item \emph{Say it:} \emph{chūnguāng}, once more
  \item \emph{Say it:} \emph{chūnguāng zuì hǎo de shíhou shù dōu kāihuā le} — When spring is at its loveliest, the trees are all in blossom.
  \item \emph{Recall:} say the Mandarin for to put up an umbrella, then the Mandarin for a thick fog, then say \emph{chūnguāng} again
\end{itemize}

\begin{tcolorbox}[breakable,colback=teal!4,colframe=teal!35!black,title={Before you move on}]
What does \zh{春光} mean? (\textquotedblleft{}The beauty of spring; spring sunshine\textquotedblright{}.) Say it once more.
\end{tcolorbox}
\section[qīngchén]{qīngchén — early in the morning; at daybreak}

\label{lesson:ZH-C519-qingchen}



\begin{quote}
Before the new one: say the Mandarin for a thick fog, then the Mandarin for the beauty of spring.
\end{quote}

\subsection*{You'll want to know: qīngchén}
\emph{qīngchén} — \textquotedblleft{}early in the morning; at daybreak\textquotedblright{}.

\begin{quote}
\emph{fùqīn měitiān qīngchén qù gōngyuán pǎobù} — My father goes running in the park early every morning.
\end{quote}

Its characters come later; for now it is a word for the ear.

\begin{cousinweb}
\emph{qīngchén} is built from \emph{qīng} 'clear' and \emph{chén} 'dawn', both inherited from Middle Chinese, \emph{tshjeng} and \emph{dzyin} (Guangyun, in Baxter's transcription), which glosses the first 'clear, still' and the second 'early' and 'bright'.
\end{cousinweb}

\begin{grammarlens}[title={What you've built}]
One more word for saying how it was.
\end{grammarlens}

\subsection*{Your turn}
\begin{itemize}
\raggedright
  \item \emph{Say it:} \emph{qīngchén}
  \item \emph{Say it:} \emph{qīngchén}, once more
  \item \emph{Say it:} \emph{fùqīn měitiān qīngchén qù gōngyuán pǎobù} — My father goes running in the park early every morning.
  \item \emph{Recall:} say the Mandarin for a thick fog, then the Mandarin for the beauty of spring, then say \emph{qīngchén} again
\end{itemize}

\begin{tcolorbox}[breakable,colback=teal!4,colframe=teal!35!black,title={Before you move on}]
What does qīngchén mean? (\textquotedblleft{}Early in the morning; at daybreak\textquotedblright{}.) Say it once more.
\end{tcolorbox}
\section[\texorpdfstring{\zh{四季如春} (sìjì rú chūn)}{四季如春 (sìjì rú chūn)}]{\zh{四季如春} (sìjì rú chūn) — mild like spring all year round (of a place)}

\label{lesson:ZH-C519-sijiruchun}



\begin{quote}
Before the new one: say the Mandarin for the beauty of spring, then the Mandarin for early in the morning.
\end{quote}

\subsection*{You'll want to know: \zh{四季如春}}
\textbf{\zh{四季如春}} — \emph{sìjì rú chūn} — \textquotedblleft{}mild like spring all year round (of a place)\textquotedblright{}.

\begin{quote}
\textbf{\zh{这儿四季如春我最喜欢这儿的阳光}} \emph{zhèr sìjì rú chūn wǒ zuì xǐhuan zhèr de yángguāng} — It's like spring here all year round, and the sunshine is what I love most.
\end{quote}

\begin{cousinweb}
\textbf{\zh{四季如春}} (\emph{sìjì rú chūn}) is built from \zh{四季} 'the four seasons', \zh{如} 'like' and \zh{春} 'spring', all inherited from Middle Chinese, \emph{sijH}, \emph{kjwijH}, \emph{nyo} and \emph{tsyhwin} (Guangyun, in Baxter's transcription), which glosses \zh{如} 'to resemble' and the last, after an older book, as the head of the four seasons: every season like the spring.
\end{cousinweb}

\begin{grammarlens}[title={What you've built}]
One more word for saying how it was.
\end{grammarlens}

\subsection*{Your turn}
\begin{itemize}
\raggedright
  \item \emph{Say it:} \emph{sìjì rú chūn}
  \item \emph{Say it:} \emph{sìjì rú chūn}, once more
  \item \emph{Say it:} \emph{zhèr sìjì rú chūn wǒ zuì xǐhuan zhèr de yángguāng} — It's like spring here all year round, and the sunshine is what I love most.
  \item \emph{Recall:} say the Mandarin for the beauty of spring, then the Mandarin for early in the morning, then say \emph{sìjì rú chūn} again
\end{itemize}

\begin{tcolorbox}[breakable,colback=teal!4,colframe=teal!35!black,title={Before you move on}]
What does \zh{四季如春} mean? (\textquotedblleft{}Mild like spring all year round (of a place)\textquotedblright{}.) Say it once more.
\end{tcolorbox}
\section[xuějǐng]{xuějǐng — a snowy landscape}

\label{lesson:ZH-C519-xuejing}



\begin{quote}
Before the new one: say the Mandarin for early in the morning, then the Mandarin for like spring all year round.
\end{quote}

\subsection*{You'll want to know: xuějǐng}
\emph{xuějǐng} — \textquotedblleft{}a snowy landscape\textquotedblright{}.

\begin{quote}
\emph{wǒ méi kàn guo zhème hǎokàn de xuějǐng} — I've never seen such a lovely snowy landscape.
\end{quote}

Its characters come later; for now it is a word for the ear.

Two third tones meet in this word: xuě is said with a rising tone before jǐng, and the spelling keeps both marks.

\begin{cousinweb}
\emph{xuějǐng} is built from \zh{雪} 'snow' and \emph{jǐng} 'a view', both inherited from Middle Chinese, \emph{sjwet} and \emph{kjaengX} (Guangyun, in Baxter's transcription), which glosses the first, after an older book, as rain that has set, and the second 'bright' and 'an image'.
\end{cousinweb}

\begin{grammarlens}[title={What you've built}]
One more word for saying how it was.
\end{grammarlens}

\subsection*{Your turn}
\begin{itemize}
\raggedright
  \item \emph{Say it:} \emph{xuějǐng}
  \item \emph{Say it:} \emph{xuějǐng}, once more
  \item \emph{Say it:} \emph{wǒ méi kàn guo zhème hǎokàn de xuějǐng} — I've never seen such a lovely snowy landscape.
  \item \emph{Recall:} say the Mandarin for early in the morning, then the Mandarin for like spring all year round, then say \emph{xuějǐng} again
\end{itemize}

\begin{tcolorbox}[breakable,colback=teal!4,colframe=teal!35!black,title={Before you move on}]
What does xuějǐng mean? (\textquotedblleft{}A snowy landscape\textquotedblright{}.) Say it once more.
\end{tcolorbox}
\section[chéngliáng]{chéngliáng — to sit in the cool (out of the sun)}

\label{lesson:ZH-C519-chengliang}



\begin{quote}
Before the new one: say the Mandarin for like spring all year round, then the Mandarin for a snowy landscape.
\end{quote}

\subsection*{You'll want to know: chéngliáng}
\emph{chéngliáng} — \textquotedblleft{}to sit in the cool (out of the sun)\textquotedblright{}.

\begin{quote}
\emph{xiàtiān de wǎnshang dàjiā dōu zài wàibian chéngliáng} — On summer evenings everyone sits outside in the cool.
\end{quote}

Its characters come later; for now it is a word for the ear.

\begin{cousinweb}
\emph{chéngliáng} is built from \emph{chéng} 'to ride on' and \emph{liáng} 'cool', both inherited from Middle Chinese, \emph{zying} and \emph{ljang} (Guangyun, in Baxter's transcription), which glosses the first 'to drive' and 'to mount' and the second 'thin' and 'cold': to take the cool air.
\end{cousinweb}

\begin{grammarlens}[title={What you've built}]
One more word for saying how it was.
\end{grammarlens}

\subsection*{Your turn}
\begin{itemize}
\raggedright
  \item \emph{Say it:} \emph{chéngliáng}
  \item \emph{Say it:} \emph{chéngliáng}, once more
  \item \emph{Say it:} \emph{xiàtiān de wǎnshang dàjiā dōu zài wàibian chéngliáng} — On summer evenings everyone sits outside in the cool.
  \item \emph{Recall:} say the Mandarin for like spring all year round, then the Mandarin for a snowy landscape, then say \emph{chéngliáng} again
\end{itemize}

\begin{tcolorbox}[breakable,colback=teal!4,colframe=teal!35!black,title={Before you move on}]
What does chéngliáng mean? (\textquotedblleft{}To sit in the cool (out of the sun)\textquotedblright{}.) Say it once more.
\end{tcolorbox}
\section[(dialogue)]{Looking back: from When the Weather Turns to Out in Nature}

\label{lesson:ZH-R519-review-514-519}



\begin{quote}
Say the Mandarin for a flash of lightning and to get caught in the rain.
\end{quote}

\subsection*{Your turn}
\begin{itemize}
\raggedright
  \item \emph{Say it:} a flash of lightning, to get caught in the rain, a thunderstorm, to be frozen stiff, to melt
  \item \emph{Say it:} bitterly cold, the four seasons, the rainy season, to escape the summer heat, to bring in the harvest
  \item \emph{Say it:} fallen leaves, a tree, a fruit tree, the woods, to ripen
  \item \emph{Say it:} to get light at dawn, light from the sun, to give off light, starlight, twilight
  \item \emph{Say it:} clear and cloudless, sunshine, solar power, to put up an umbrella, a thick fog
  \item \emph{Say it:} the beauty of spring, early in the morning, like spring all year round, a snowy landscape, to sit in the cool
  \item \emph{Read it:} \zh{季} \zh{树} \zh{光} \zh{阳} — each one aloud, then a word that holds it
\end{itemize}

\begin{tcolorbox}[breakable,colback=teal!4,colframe=teal!35!black,title={Before you move on}]
A flash of lightning? (\emph{shǎndiàn}.) To get caught in the rain? (\emph{línyǔ}.) A thunderstorm? (\emph{léiyǔ}.) To be frozen stiff? (\emph{dòngjiāng}.) To melt? (\emph{rónghuà}.) Bitterly cold? (\emph{hánlěng}.) The four seasons? (\textbf{\zh{四季}}, \emph{sìjì}.) The rainy season? (\textbf{\zh{雨季}}, \emph{yǔjì}.) To escape the summer heat? (\emph{bìshǔ}.) To bring in the harvest? (\emph{shōuhuò}.) Fallen leaves? (\emph{luòyè}.) A tree? (\textbf{\zh{树}}, \emph{shù}.) A fruit tree? (\textbf{\zh{果树}}, \emph{guǒshù}.) The woods? (\emph{shùlín}.) To ripen? (\emph{chéngshú}.) To get light at dawn? (\emph{tiānliàng}.) Light from the sun? (\textbf{\zh{光}}, \emph{guāng}.) To give off light? (\textbf{\zh{发光}}, \emph{fāguāng}.) Starlight? (\textbf{\zh{星光}}, \emph{xīngguāng}.) Twilight? (\emph{huánghūn}.) Clear and cloudless? (\emph{qínglǎng}.) Sunshine? (\textbf{\zh{阳光}}, \emph{yángguāng}.) Solar power? (\textbf{\zh{太阳能}}, \emph{tàiyángnéng}.) To put up an umbrella? (\emph{dǎsǎn}.) A thick fog? (\emph{dàwù}.) The beauty of spring? (\textbf{\zh{春光}}, \emph{chūnguāng}.) Early in the morning? (\emph{qīngchén}.) Like spring all year round? (\textbf{\zh{四季如春}}, \emph{sìjì rú chūn}.) A snowy landscape? (\emph{xuějǐng}.) To sit in the cool? (\emph{chéngliáng}.)
\end{tcolorbox}
